% Reference sheet — the module's results and notation at their densest, one artifact per topic or
% per exam. Formulas and conditions only; anything that argues belongs in revision notes.
% Copy it into the folder it serves and replace every <...> placeholder. The input line below
% resolves the preamble from the templates folder or from a unit folder two levels down.
\documentclass[a4paper,10pt]{article}
\IfFileExists{preamble.tex}{\input{preamble.tex}}{\input{../../templates/preamble.tex}}
\begin{document}
\ArtifactHeader{\ModuleCode}{<Source-map key, or the exam this serves>}{<what this sheet covers>}%
  {<YYYY-MM-DD>}{<the units and chapters the entries were taken from>}

\section{Notation}
\begin{itemize}
  \item <symbol> --- <what it denotes in this module, where the module has a house convention
        that differs from the textbook's.>
\end{itemize}

\section{<A group of results>}
\begin{itemize}
  \item \textbf{<name>.} <the statement> \quad <the condition it needs, always carried with it:
        a formula written without its hypothesis is the error this sheet exists to stop.>
  \item \textbf{<name>.}
        \[
          \text{<a result that needs a display>}
        \]
        <its condition>
\end{itemize}

\section{<A second group>}
\begin{itemize}
  \item \textbf{<name>.} <the statement> \quad <its condition>
\end{itemize}

\begin{warning}
<The pair of results that get used in each other's place, and the one word that tells them apart.>
\end{warning}
\end{document}
